\PassOptionsToPackage{dvipsnames,svgnames,table}{xcolor}
\documentclass[10pt]{article}
\usepackage[a4paper,margin=22mm]{geometry}
\usepackage[no-math]{fontspec}\setmainfont[SmallCapsFont={Latin Modern Roman Caps}]{Latin Modern Roman}
\usepackage{amsmath,amssymb,amsthm,mathtools,graphicx,xcolor,hyperref,xurl,xspace,ifthen,enumerate,textcomp}
\usepackage[shortlabels]{enumitem}
\setlength{\emergencystretch}{4em}
\newtheorem{theorem}{Theorem}[section]
\newtheorem{theo}[theorem]{Theorem}
\newtheorem{thm}[theorem]{Theorem}
\newtheorem{lemma}[theorem]{Lemma}
\newtheorem{lemm}[theorem]{Lemma}
\newtheorem{prop}[theorem]{Proposition}
\newtheorem{coro}[theorem]{Corollary}
\newtheorem{corollary}[theorem]{Corollary}
\newtheorem{fact}[theorem]{Fact}
\newtheorem{claim}[theorem]{Claim}
\theoremstyle{definition}
\newtheorem{defi}[theorem]{Definition}
\newtheorem{definition}[theorem]{Definition}
\newtheorem{rema}[theorem]{Remark}
\newtheorem{remark}[theorem]{Remark}
\newtheorem{exam}[theorem]{Example}
\newtheorem{example}[theorem]{Example}
\newenvironment{enonce*}[2][]{\par\medskip\noindent\textbf{#2.}\itshape}{\par\medskip}
\providecommand{\eg}{e.g.\ }
\providecommand{\ie}{i.e.\ }
\providecommand{\resp}{resp.\ }
\providecommand{\cf}{cf.\ }
\providecommand{\longthanks}[1]{\par\smallskip\noindent #1\par\smallskip}
\newtheorem{proposition}[theorem]{Proposition}
\newtheorem{conjecture}[theorem]{Conjecture}
\numberwithin{equation}{section}
\newcommand{\R}{{\mathbb R}}
\newcommand{\Z}{{\mathbb Z}}
\newcommand{\N}{{\mathbb N}}
\newcommand{\expref}[2]{\texorpdfstring{\hyperref[#2]{#1~\ref{#2}}}{#1~\ref{#2}}}
\newcommand{\epfmtdoi}[1]{\href{https://doi.org/#1}{doi:#1}}
\newcommand{\epfmt}[2]{#1:\,#2}
\newcommand{\cclass}[1]{\textsf{#1}}
\newcommand{\opt}{\textsc{OPT}}
\newcommand{\etal}{et al.\ }
\newcommand{\eps}{\epsilon}
\newcommand{\st}{S} %
\newcommand{\Algorithm}[1]{{\textrm{#1}}} %
\newcommand{\sbg}{\Algorithm{SSF-W}} %
\newcommand{\sug}{\Algorithm{SSF}} %
\newcommand{\mmug}{\Algorithm{SSF-ID}} %
\newcommand{\lf}{\Algorithm{LF}}
\newcommand{\ld}{\Algorithm{LD}}
\newcommand{\bwf}{\Algorithm{BWF}}
\newcommand{\lwf}{\Algorithm{LWF}}
\newcommand{\fifo}{\Algorithm{FIFO}}
\newcommand{\partially}{\mathcal{A}} %
\newcommand{\zeroly}{\mathcal{B}}      %
\newcommand{\bwfw}{\Algorithm{BWF-W}} %
\newcommand{\srfw}{\Algorithm{SRF-W}} %
\newcommand{\srf}{\Algorithm{SRF}} %
\begin{document}
\begingroup\raggedright
\section*{1599. Scalable immediate-dispatch nonmigratory scheduling for maximum delay factor}
Chandra Chekuri, Sungjin Im, and Benjamin Moseley. \emph{Online Scheduling to Minimize Maximum Response Time and Maximum Delay Factor}. Theory of Computing volume 8 article 7 (2012), official complete journal PDF acquired 2026-09-30; canonical DOI 10.4086/toc.2012.v008a007 verified against official landing; printed header identity cross-checked and any variant preserved. \url{https://theoryofcomputing.org/articles/v008a007/}. CC BY3.0.
\paragraph{Selected problem and scope (editorial).} Exact original Theorem2.5 only: SSF-ID is a (1+epsilon)-speed O(1/epsilon)-competitive online algorithm for maximum delay factor on m identical machines. Complete online request/slack/deadline/delay-factor and competitive-analysis definitions, explicit zero-cost preemption, epsilon in (0,1], speed-augmentation and S\_i>=ell\_i conventions, unicast processing model, full SSF definition with tie-breaking/preemption and its source-supplied single-machine proof, immediate cumulative-class-volume dispatch algorithm, assigned/processed/residual volume invariants and optimal-backlog corollary, and entire unwrapped critical-interval competitive proof are retained. Ordinary online competitive analysis, finite workload accounting, dyadic sums and elementary inequalities are prerequisites; every new dispatch/invariant/critical-interval argument is supplied. Native inactive iffalse alternative SSF rule at604–609 is excluded with all physical newlines preserved; it was never active author text. Single-machine and helper invariants are uncounted support; independent broadcast1600 and weighted/nonuniform-page variants are unselected.
\paragraph{Difficulty (editorial).} The target must dispatch by cumulative assigned class volume, rather than current unfinished queue size, and prove separate assigned-volume, processed-volume and residual-volume balance invariants before comparison with a potentially migratory optimum is available. The complete authored busy-interval closing at804–838 converts those local bounds into a machine-count-independent backlog and competitive guarantee. This is one H2 online dispatch/structural-balance task above a direct single-machine shortest-slack workload calculation.
\paragraph{Formatting (editorial).} Exact original human mathematical body and order retained; standalone typography, counter restoration and declared noncore printed locators only. Original comment prose excluded with percent guards retained; source typography and counters restored without upstream class execution. No mathematical repair or reconstruction.
\paragraph{Original source quirks (editorial).} The source single-machine introduction simultaneously mentions nonuniform sizes and then says uniform sizes at582–588; its full original model/algorithm and proof are retained without harmonization. The primary backlog display at825–826 omits the subtraction for optimal processing explicitly mentioned immediately beforehand at821–822; display and prose are preserved together. Source class-index conventions, alpha=max-delay versus ratio wording, and original max\{16,2alpha*/epsilon\} closing are unchanged; no mathematical repair or refereeing is supplied.
\endgroup
\clearpage
\setcounter{section}{1}
We assume that requests arrive \emph{online}. A request $J_i$ and its
properties are known only when it arrives to the system.  We use
$a_i$ to denote its arrival time and $\ell_i$ to denote its
processing time. In the weighted case $J_i$ has a non-negative
weight $w_i$; the unweighted case corresponds to $w_i=1$ for all
$i$. Consider an online scheduling algorithm $A$. Let $f^A_i$ denote
the completion time or finish time of $J_i$ under $A$.  The response
time of $J_i$ under $A$ is $f^A_i - a_i$, in other words the total
duration spent by $J_i$ in the system before its processing is
finished. Now we define the delay factor metric. Here it is
assumed that each request $J_i$ has a \emph{deadline} $d_i$ that is
known upon its arrival. We refer to the quantity $S_i = (d_i-a_i)$
as the \emph{slack} of request $J_i$. If $f^A_i$ is the finish time
of $J_i$ under some schedule $A$ then its delay factor is defined as
$\max \{ 1, {(f^A_i - a_i)}/{(d_i - a_i)}\}$. Thus, delay factor
measures the ratio of the response time and the slack, unless the
request finishes by its deadline in which case we set its delay
factor to $1$.  In this paper we are interested in algorithms that
minimize the maximum response time and the maximum delay factor.
Given an online request sequence $\sigma$ and an algorithm $A$, let
$\alpha^A(\sigma) = \max_{J_i \in \sigma} \gamma^A_i$ where
$\gamma^A_i$ is either the response time or the delay factor of
$J_i$ in $A$; in the weighted case $\alpha^A(\sigma) = \max_{J_i \in
\sigma} w_i \gamma^A_i$.  We measure the performance of an algorithm
$A$ via worst-case competitive analysis: $A$ is $r$-competitive if
for all request sequences $\sigma$ we have $\alpha^A(\sigma) \le r
\alpha^*(\sigma)$ where $\alpha^*(\sigma)$ is the value of an
optimum offline schedule for $\sigma$.

\subsection{Results}
We give the first non-trivial results for online scheduling to minimize the
maximum delay factor and weighted response time in both the unicast
and broadcast settings. Throughout we assume that requests are
allowed to be preempted without any penalty.  We remark that
weighted response time and delay factor, though formally not
equivalent, behave similarly in terms of algorithmic development. At
a heuristic level, one can interpret the term ${1}/{(d_i-a_i)}$ in
the delay factor metric as the weight $w_i$. For this reason, we
mainly discuss results for delay factor below and point out how they
generalize to weighted response time.

We resort to resource augmentation analysis, introduced by of
Kalyanasundaram and Pruhs~\cite{KalyanasundaramP95}, to overcome the
above lower bounds. In this analysis the online algorithm is given a
faster machine than the optimal offline algorithm. For $s \ge 1$, an
algorithm $A$ is said to be $s$-speed $r$-competitive if $A$, when
given $s$-speed machine(s), achieves a competitive ratio of $r$
compared to the optimum  schedule given 1-speed machine(s).
We prove the following.
\begin{itemize}
\item For unicast setting, for any $\eps \in (0,1]$, there are
  $(1+\eps)$-speed $O(1/\eps)$-competitive algorithms in both single
  and multiple machine cases. Moreover, the algorithm for the multiple
  machine case immediately dispatches an arriving request to a
  machine, and is non-migratory. An algorithm is non-migratory
  if it processes each request on a single machine to which it is
  first assigned.

\item For broadcast setting, for any $\eps \in (0,1]$, there is a
  $(1+\eps)$-speed $O(1/\eps^2)$-competitive algorithm for uniform
  pages.  For non-uniform sized pages, for any $\eps \in (0,1]$, there
  is a $(1+\eps)$-speed $O(1/\eps^3)$-competitive algorithm.
\end{itemize}

\paragraph{Notation}
We let $S_i = d_i - a_i$ denote the slack of $J_i$ in the unicast
setting. When requests have non-uniform processing times (or lengths) we
use $\ell_i$ to denote the length of $J_i$. We assume without loss
of generality that $S_i \ge \ell_i$. In the broadcast setting,
$J_{p,i}$ denotes the $i$'th request for page $p$. We assume that the
requests for a page are ordered by arrival time and hence $a_{p,i}
\le a_{p,j}$ for $i < j$. In both settings we use $\Delta$ to denote
the ratio of maximum slack to the minimum slack in a given request
sequence. When pages have non-uniform sizes, $\ell_p$ denotes the size of page $p$.  For an interval $I = [a,b]$ on the real line we use $|I|$
for the length of the interval $(b-a)$. We say a request is alive
 at $t$ if has arrived by $t$ but has not yet finished.
\setcounter{section}{1}
\section{Unicast scheduling}
\label{sec:unicast} In this section we address the unicast setting
where requests are independent and consider the delay factor metric.
For a request $J_i$, recall that $a_i,d_i,\ell_i,f_i$ denote the
arrival time, deadline, processing time/size, and finish time,
respectively. An instance with all $\ell_i = 1$ (or more generally
the processing times are uniform) is referred to as a uniform
processing time instance. It is easy to see that preemption does not
help much for uniform processing time instances when the maximum
delay factor metric is considered. Assuming that the processing
times are integer valued, in the single machine setting one can
reduce an instance with non-uniform processing times to an instance
with uniform processing times as follows:  replace $J_i$, with
processing time $\ell_i$ by $\ell_i$ uniform sized requests with
the same arrival time and deadline as that of $J_i$.
\setcounter{subsection}{0}\setcounter{theorem}{2}
\subsection{Single machine scheduling}
\label{unicast-single}

In this section we consider a simple greedy algorithm for minimizing
the maximum delay factor on a single machine when requests have
non-uniform sizes.  We analyze the simple shortest-slack-first
($\sug$) algorithm which at any time $t$ schedules the request with
the shortest slack.  Recall that the slack of a request $J_i$ is $d_i
- a_i$ and that we have assumed without loss of generality that all
requests have uniform sizes.

\begin{center}
\begin{tabular}[r]{|c|}
\hline
\textbf{Algorithm}: \sug \\

\\

\begin{minipage}{13cm}
\begin{itemize}
\item Let $Q(t)$ be the set of alive requests at $t$.
\item Let $J_i$ be the request with the minimum slack among requests
  in $Q(t)$, ties broken by arrival time.
\item Preempt the current request
  and schedule $J_i$ if it is not being processed.






\end{itemize}
\end{minipage}\\\\

\hline
\end{tabular}
\end{center}


\begin{theorem}
\label{thm:ssf-single-machine} The algorithm $\sug$ is an $(1 +
\eps)$-speed ${1/\eps}$-competitive online algorithm for minimizing the maximum delay factor in the unicast setting.
\end{theorem}
\begin{proof}
  Consider an arbitrary request sequence $\sigma$ and let $\alpha^\sug$ be
  the maximum delay factor achieved by $\sug$ on $\sigma$. If $\alpha^\sug
  = 1$ there is nothing to prove, so assume that $\alpha^\sug > 1$.  Let
  $J_i$ be the request that witnesses $\alpha^\sug$, that is $\alpha^\sug = (f_i
  - a_i)/S_i$. Note that $\sug$ does not process any request with slack
  more than $S_i$ in the interval $[a_i, f_i]$. Let $t$ be the minimum
  value less than or equal to $a_i$ such that $\sug$ was busy processing only
  requests with slack at most $S_i$ in the interval $[t, f_i]$. It
  follows that $\sug$ had no requests with slack $\le S_i$ just before
  $t$. The total work that $\sug$ processed in $[t,f_i]$ on requests
  with slack less than equal to $S_i$ is $(1+\eps)(f_i-t)$ and all
  these requests arrive in the interval $[t,f_i]$. An optimal offline
  algorithm with $1$-speed can do total work of at most $(f_i-t)$ in
  the interval $[t,f_i]$ and hence the earliest time by which it can
  finish these requests is $f_i + \eps (f_i-t) \ge f_i + \eps(f_i -
  a_i)$. Since all these requests have slack at most $S_i$ and have
  arrived before $f_i$, it follows that $\alpha^* \ge \eps(f_i -
  a_i)/S_i$ where $\alpha^*$ is the maximum delay factor of the
  optimal offline algorithm with $1$-speed machine. Therefore, we have
  that $\alpha^{\sug}/\alpha^* \le 1/\eps$.
\end{proof}

\begin{remark}
  For uniform processing time requests, the algorithm that \emph{non-preemptively} schedules requests with the shortest slack is a
  $(1+\eps)$-speed $2/\eps$-competitive online  algorithm for
  minimizing the maximum delay factor.
\end{remark}

\subsection{Multiple machine scheduling}
\label{unicast-multi}

We now consider minimizing the maximum delay factor when there are $m$
machines. In the multiple machine setting, at any time a machine can
choose to process at most one request and a request can be processed
by at most one machine at a given time. A request is allowed to
migrate and be processed by different machines at different times; as
we remarked earlier, a schedule or algorithm is non-migratory if it
does not migrate any request. To adapt $\sug$ to this setting we take
intuition from previous work on minimizing $L_k$ norms of flow time
and stretch~\cite{BansalP03,AvrahamiA03,ChekuriGKK04}. We develop an
algorithm that immediately dispatches an arriving request to a
machine, and further does not migrate an assigned request to a
different machine once it is assigned. Each machine essentially runs
the single machine $\sug$ algorithm and thus the only remaining
ingredient to describe is the dispatching rule. For this purpose the
algorithm groups requests into classes based on their slack. A request
$J_i$ is said to be in class $k$ if $S_i \in [2^k, 2^{k+1})$. The
algorithm maintains the total processing time of requests (referred to
as \emph{volume}) that have been assigned to machine $x$ in each class
$k$. Let $U^x_{=k}(t)$ denote the total processing time of requests
assigned to machine $x$ by time $t$ of class $k$. With this notation,
the algorithm $\mmug$ (for $\sug$ with immediate dispatch) can be
described.

\begin{center}
\begin{tabular}[r]{|c|}
\hline
\textbf{Algorithm}: \mmug \\

\\

\begin{minipage}{13cm}
\begin{itemize}
\item When a new request $J_i$ of class $k$ arrives at time $t$,
assign it to a machine $x$ where $U^x_{=k}(t) = \min_y U^y_{=k}(t)$.
\item Use $\sug$ on each machine separately.
\end{itemize}
\end{minipage}\\\\

\hline
\end{tabular}
\end{center}

The rest of this section is devoted to the proof of the following
theorem.
\begin{theorem}
\label{thm:multiple} $\mmug$ is a $(1+ \eps)$-speed $O({1/
\eps})$-competitive online algorithm for minimizing the maximum delay
factor on $m$ identical machines.
\end{theorem}

We need a fair amount of notation. For each time $t$, machine $x$, and
class $k$ we define several quantities. For example $U^x_{=k}(t)$ is
the total volume assigned to machine $x$ in class $k$ by time $t$. We
use the predicate ``$\le k$'' to indicate classes $1$ to $k$.  Thus $U^x_{\le
  k}(t)$ is the total volume assigned to machine $x$ in classes $1$ to
$k$. We let $R^x_{=k}(t)$ to denote the remaining processing time on
machine $x$ at time $t$ and let $P^x_{=k}(t)$ denote the total volume
that $x$ has finished on requests in class $k$ by time $t$. Note that
$P^x_{=k}(t) = U^x_{=k}(t) - R^x_{=k}(t)$. All these quantities refer
to the algorithm $\mmug$. We use $V^*_{=k}(t)$ and $V_{=k}(t)$ to
denote the remaining volume of requests in class $k$ in an optimal
offline algorithm with speed $1$ and $\mmug$ with speed $(1+\eps)$,
respectively. Observe that $V_{=k}(t) = \sum_x R^x_{=k}(t)$.  The
quantities $V^*_{\le k}(t)$ and $V_{\le k}(t)$ are defined
analogously.

The algorithm $\mmug$ balances the amount of processing time for
requests with similar slack.  Note that the assignment of requests is
not based on the current volume of unfinished requests on the
machines, rather the assignment is based on the volume of requests
that were assigned in the past to different machines.  We begin our
proof by showing that for any $k$ the total volume of requests of
class at most $k$ that is assigned is balanced across the machines.
This easily follows from the dispatching policy of the
algorithm. Several of the lemmas below and their proofs
follow from the work in~\cite{AvrahamiA03,ChekuriGKK04}.

\begin{proposition}
\label{observation} 
For any time $t$ and two machines $x$ and $y$,
$|U^{x}_{=k}(t) - U^{y}_{=k}(t)| \leq 2^{k+1}$.  This also implies
that $|U^{x}_{\leq k}(t) - U^{y}_{\leq k}(t)| \leq 2^{k+2}$.
\end{proposition}
\begin{proof}
  The first inequality holds since all of the requests of class $k$
  are of size $\leq 2^{k+1}$.  The second inequality follows easily
  from the first.
\end{proof}




\begin{lemma}
  \label{lem:processbound} Consider any two machines $x$ and $y$.  At any time $t$, the
  difference in volume of requests that already have been
  processed is bounded as $|P^{x}_{\leq k} (t) - P^{y}_{\leq k} (t)|
  \leq 2^{k+2}$.
\end{lemma}

\begin{proof}
  Suppose the lemma is false. Then there is the first time $t_0$ when
  $P^{x}_{\leq k} (t_0) - P^{y}_{\leq k} (t_0) = 2^{k+2}$ and small
  constant $\delta t > 0$ such that $P^{x}_{\leq k} (t_0 + \delta t) -
  P^{y}_{\leq k} (t_0 + \delta t) > 2^{k+2}$. Let $t' = t_0+\delta t$.
  For this to occur, $x$ processes a request of class $\le k$ during
  the interval $I = [t_0, t']$ while $y$ processes a request of class
  $> k$. Since each machine uses $\sug$, it must be that $y$ had no
  requests in classes $\le k$ during $I$ which implies that $U^y_{\le
    k}(t') = P^y_{\le k}(t')$. Therefore,
\[
U^{y}_{\leq k} (t') = P^{y}_{\leq k} (t') < P^{x}_{\leq k} (t') - 2^{k+2}
\le U^{x}_{\leq k} (t') - 2^{k+2}\,,
\]
since $P^{x}_{\leq k} (t') \le U^{x}_{\leq k} (t')$. However, this implies
that
\[
U^{y}_{\leq k} (t') < U^{x}_{\leq k} (t')  - 2^{k+2}\,,
\]
a contradiction to \expref{Proposition}{observation}.
\end{proof}


\begin{lemma}
  \label{lem:closebound} At any time $t$ the difference between the
  residual volume of requests that needs to be processed, on any two
  different machines $x$ and $y$, is bounded as $|R^{x}_{\leq k} (t) -
  R^{y}_{\leq k}(t)| \leq 2^{k+3}$.
\end{lemma}

\begin{proof}
Combining \expref{Proposition}{observation}, \expref{Lemma}{lem:processbound},
and the fact that $R^{x}_{\leq k}(t) = U^{x}_{\leq k}(t) - P^{x}_{\leq k}(t)$ for any $x$ and $k$ by definition then,
\[
|R^{x}_{\leq k} (t) -  R^{y}_{\leq k}(t)| \le
|U^{x}_{\leq k} (t) -  U^{y}_{\leq k}(t)| + |P^{x}_{\leq k} (t) -  P^{y}_{\leq k}(t)| \le 2^{k+3}\,.
\]
\end{proof}


\begin{corollary}
  \label{cor:optarrival} At any time $t$, $V^*_{\leq k} (t) \geq V_{\leq
    k} (t) - m2^{k+3}$.
\end{corollary}

%
%
%
%
%

We are now ready to upper bound the competitiveness of \mmug~ when given $(1+\eps)$-speed in a similar fashion to the single machine
case. Consider an arbitrary request sequence $\sigma$ and let $J_i$ be
the request that witnesses the maximum delay factor $\alpha^\mmug$ of $\mmug$. Let $k$ be the class of $J_i$ and let $x$ be the machine $J_i$ was processed on by $\mmug$.  We know that $\alpha^\mmug = (f_i - a_i)/S_i$ by definition of $J_i$. We use $\alpha^*$ to denote the delay factor of some fixed optimal
algorithm that uses $m$ machines of speed $1$.

Let $t$ be the last time before $a_i$ when machine $x$ processed a request
of class $> k$ under $\mmug$'s schedule. Note that $t \leq a_i$ since $x$ does not process any
request of class $> k$ in the interval $[a_i,f_i]$. At time $t$ we
know by \expref{Corollary}{cor:optarrival} that $V^*_{\leq k} (t) \geq
V_{\leq k} (t) - m2^{k+3}$. If $f_i \le a_i + 2^{k+4}$ then $\mmug$
achieves a competitive ratio of $16$ since $J_i$ is in class
$k$. Thus we will assume from now on that $f_i > a_i + 2^{k+4}$.

During the interval $I = [t, f_i)$, $\mmug$ completes a total volume of
$(1 + \eps)(f_i - t)$ work on machine $x$. Using \expref{Lemma}{lem:processbound}, any
other machine $y$ also processes a volume of $(1+\eps)(f_i-t) -
2^{k+3}$ work during $I$ of requests in classes at most $k$ . Thus the total volume processed by $\mmug$ during $I$ of requests in classes at most $ k$ is at least $m(1+\eps)(f_i-t) -
m2^{k+3}$.  During $I$, the optimal algorithm schedules at most a
$m(f_i-t)$ volume of work for requests in classes at most $k$.  Combining this with
\expref{Corollary}{cor:optarrival}, we see that
\begin{eqnarray*}
V^*_{\le k}(f_i) & \ge & V_{\le k}(t) - m2^{k+3} + m(1+\eps)(f_i-t)  - m2^{k+3} \\
  & \ge & V_{\le k}(t) + m(1+\eps)(f_i-t) - m2^{k+4} \ge   \eps m (f_i - t)\,.
\end{eqnarray*}
In the last inequality we use the fact that $f_i-t \ge f_i -
a_i \ge 2^{k+4}$. Without loss of generality assume that no
requests arrives exactly at $f_i$. Therefore $V^*_{\le k}(f_i)$ is the
total volume of requests in classes $1$ to $k$ that the optimal
algorithm has left to finish at time $f_i$ and all these requests have
arrived before $f_i$. The earliest time that the optimal algorithm can
finish all these requests is $f_i + \eps (f_i-t)$ and therefore it
follows that $\alpha^* \ge \eps (f_i - t)/2^{k+1}$. Since $\alpha^\mmug \le (f_i-a_i)/2^k$ and $t \leq a_i$, it follows that $\alpha^\mmug \le 2 \alpha^*/\eps$.

Thus $\alpha^\mmug \le \max\{16, 2\alpha^*/\eps\}$ which completes the proof
of \expref{Theorem}{thm:multiple}.
\clearpage
\begingroup\raggedright
%
%
%
%
%
\providecommand{\bibhead}[1]{}
%
\expandafter\ifx\csname pdfbookmark\endcsname\relax%
  \providecommand{\tocrefpdfbookmark}{}
\else\providecommand{\tocrefpdfbookmark}{%
   \hypertarget{tocreferences}{}%
   \pdfbookmark[1]{References}{tocreferences}}%
\fi

\tocrefpdfbookmark
\begin{thebibliography}{10}

\bibitem{AcharyaFZ95}\bibhead{AcharyaFZ95}
{\sc Swarup Acharya, Michael Franklin, and Stanley Zdonik}: Dissemination-based
  data delivery using broadcast disks.
\newblock {\em IEEE Personal Communications}, 2(6):50--60, Dec 1995.
\newblock [\epfmtdoi{10.1109/98.475988}]

\bibitem{AksoyF98}\bibhead{AksoyF98}
{\sc Demet Aksoy and Michael Franklin}: {R x W}: A scheduling approach for
  large-scale on-demand data broadcast.
\newblock {\em IEEE/ACM Transactions on Networking}, 7(6):846--860, 1999.
\newblock [\epfmtdoi{10.1109/90.811450}]

\bibitem{AvrahamiA03}\bibhead{AvrahamiA03}
{\sc Nir Avrahami and Yossi Azar}: Minimizing total flow time and total
  completion time with immediate dispatching.
\newblock {\em Algorithmica}, 47:253--268, 2007.
\newblock Preliminary version in
  \href{http://dx.doi.org/10.1145/777412.777415}{SPAA'03}.
\newblock [\epfmtdoi{10.1007/s00453-006-0193-6}]

\bibitem{BansalCKN05}\bibhead{BansalCKN05}
{\sc Nikhil Bansal, Moses Charikar, Sanjeev Khanna, and Joseph~(Seffi) Naor}:
  Approximating the average response time in broadcast scheduling.
\newblock In {\em Proc. 16th Ann. ACM-SIAM Symp. on Discrete Algorithms
  (SODA'05)}, pp. 215--221. SIAM, 2005.

\bibitem{BansalCS08}\bibhead{BansalCS08}
{\sc Nikhil Bansal, Don Coppersmith, and Maxim Sviridenko}: Improved
  approximation algorithms for broadcast scheduling.
\newblock {\em SIAM J. Comput.}, 38(3):1157--1174, 2008.
\newblock Preliminary version in
  \href{http://dx.doi.org/10.1145/1109557.1109596}{SODA'06}.
\newblock [\epfmtdoi{10.1137/060674417}]

\bibitem{BansalKN10}\bibhead{BansalKN10}
{\sc Nikhil Bansal, Ravishankar Krishnaswamy, and Viswanath Nagarajan}: Better
  scalable algorithms for broadcast scheduling.
\newblock In {\em Proc. 37th Internat. Colloq. on Automata, Languages and
  Programming (ICALP'10)}, pp. 324--335. Springer, 2010.
\newblock [\epfmtdoi{10.1007/978-3-642-14165-2\_28}]

\bibitem{BansalP03}\bibhead{BansalP03}
{\sc Nikhil Bansal and Kirk Pruhs}: Server scheduling in the {L$_{\mbox{p}}$}
  norm: A rising tide lifts all boat.
\newblock In {\em Proc. 35th STOC}, pp. 242--250. ACM Press, 2003.
\newblock [\epfmtdoi{10.1145/780542.780580}]

\bibitem{Bar-NoyBNS02}\bibhead{Bar-NoyBNS02}
{\sc Amotz Bar-Noy, Randeep Bhatia, Joseph~(Seffi) Naor, and Baruch Schieber}:
  Minimizing service and operation costs of periodic scheduling.
\newblock {\em Math. Oper. Res.}, 27(3):518--544, 2002.
\newblock [\epfmtdoi{10.1287/moor.27.3.518.314}]

\bibitem{Bar-NoyGKNSS09}\bibhead{Bar-NoyGKNSS09}
{\sc Amotz Bar-Noy, Sudipto Guha, Yoav Katz, Joseph~(Seffi) Naor, Baruch
  Schieber, and Hadas Shachnai}: Throughput maximization of real-time
  scheduling with batching.
\newblock {\em ACM Transactions on Algorithms}, 5(2):18:1--18:17, 2009.
\newblock [\epfmtdoi{10.1145/1497290.1497294}]

\bibitem{BartalM00}\bibhead{BartalM00}
{\sc Yair Bartal and S.~Muthukrishnan}: Minimizing maximum response time in
  scheduling broadcasts.
\newblock In {\em Proc. 11th Ann. ACM-SIAM Symp. on Discrete Algorithms
  (SODA'00)}, pp. 558--559. SIAM, 2000.
\newblock [\epfmt{acm}{338605}]

\bibitem{BenderCM98}\bibhead{BenderCM98}
{\sc Michael~A. Bender, Soumen Chakrabarti, and S.~Muthukrishnan}: Flow and
  stretch metrics for scheduling continuous job streams.
\newblock In {\em Proc. 9th Ann. ACM-SIAM Symp. on Discrete Algorithms
  (SODA'98)}, pp. 270--279. SIAM, 1998.
\newblock [\epfmt{acm}{314613.314715}]

\bibitem{BenderCT08}\bibhead{BenderCT08}
{\sc Michael~A. Bender, Rapha{\"e}l Clifford, and Kostas Tsichlas}: Scheduling
  algorithms for procrastinators.
\newblock {\em J. Scheduling}, 11(2):95--104, 2008.
\newblock [\epfmtdoi{10.1007/s10951-007-0038-4}]

\bibitem{BenderMR04}\bibhead{BenderMR04}
{\sc Michael~A. Bender, S.~Muthukrishnan, and Rajmohan Rajaraman}:
  Approximation algorithms for average stretch scheduling.
\newblock {\em J. Scheduling}, 7(3):195--222, 2004.
\newblock [\epfmtdoi{10.1023/B:JOSH.0000019681.52701.8b}]

\bibitem{BoggiaCMM05}\bibhead{BoggiaCMM05}
{\sc Gennaro Boggia, Pietro Camarda, Luigi Mazzeo, and Marina Mongiello}:
  Performance of batching schemes for multimedia-on-demand services.
\newblock {\em IEEE Transactions on Multimedia}, 7(5):920--931, 2005.
\newblock [\epfmtdoi{10.1109/TMM.2005.854383}]

\bibitem{Burnsb08}\bibhead{Burnsb08}
{\sc Alan Burns and Sanjoy Baruah}: Sustainability in real-time scheduling.
\newblock {\em Journal of Computing Science and Engineering}, 2(1):74--97,
  2008.
\newblock \url{http://jcse.kiise.org/PublishedPaper/topic_abstract.asp?idx=15}.

\bibitem{ChanLTW04}\bibhead{ChanLTW04}
{\sc Wun-Tat Chan, Tak~Wah Lam, Hing-Fung Ting, and Prudence W.~H. Wong}: New
  results on on-demand broadcasting with deadline via job scheduling with
  cancellation.
\newblock In {\em Proc. 10th Ann. Internat. Computing and Combinatorics Conf.
  (COCOON'04)}, pp. 210--218. Springer Berlin / Heidelberg, 2004.
\newblock [\epfmtdoi{10.1007/978-3-540-27798-9\_24}]

\bibitem{ChangEGK08}\bibhead{ChangEGK08}
{\sc Jessica Chang, Thomas Erlebach, Renars Gailis, and Samir Khuller}:
  Broadcast scheduling: Algorithms and complexity.
\newblock {\em ACM Transactions on Algorithms}, 7(4):47:1--47:14, 2011.
\newblock Preliminary version in
  \href{https://dl.acm.org/citation.cfm?doid=1347082.1347134}{SODA'08}.
\newblock [\epfmtdoi{10.1145/2000807.2000815}]

\bibitem{CharikarK06}\bibhead{CharikarK06}
{\sc Moses Charikar and Samir Khuller}: A robust maximum completion time
  measure for scheduling.
\newblock In {\em Proc. 17th Ann. ACM-SIAM Symp. on Discrete Algorithms
  (SODA'06)}, pp. 324--333. SIAM, 2006.
\newblock [\epfmtdoi{10.1145/1109557.1109594}]

\bibitem{ChekuriGIKLMMR10}\bibhead{ChekuriGIKLMMR10}
{\sc Chandra Chekuri, Avigdor Gal, Sungjin Im, Samir Khuller, Jian Li, Richard
  McCutchen, Benjamin Moseley, and Louiqa Raschid}: New models and algorithms
  for throughput maximization in broadcast scheduling - (extended abstract).
\newblock In {\em Proc. 8th Workshop on Approximation and Online Algorithms
  (WAOA'10)}, pp. 71--82. Springer, 2010.
\newblock [\epfmtdoi{10.1007/978-3-642-18318-8\_7}]

\bibitem{ChekuriGKK04}\bibhead{ChekuriGKK04}
{\sc Chandra Chekuri, Ashish Goel, Sanjeev Khanna, and Amit Kumar}:
  Multi-processor scheduling to minimize flow time with $\epsilon$ resource
  augmentation.
\newblock In {\em Proc. 36th STOC}, pp. 363--372. ACM Press, 2004.
\newblock [\epfmtdoi{10.1145/1007352.1007411}]

\bibitem{ChekuriIM09lwf}\bibhead{ChekuriIM09lwf}
{\sc Chandra Chekuri, Sungjin Im, and Benjamin Moseley}: Longest wait first for
  broadcast scheduling [extended abstract].
\newblock In {\em Proc. 7th Workshop on Approximation and Online Algorithms
  (WAOA'09)}, pp. 62--74. Springer, 2009.
\newblock [\epfmtdoi{10.1007/978-3-642-12450-1\_6}]

\bibitem{ChekuriIM09esa}\bibhead{ChekuriIM09esa}
{\sc Chandra Chekuri, Sungjin Im, and Benjamin Moseley}: Minimizing maximum
  response time and delay factor in broadcast scheduling.
\newblock In {\em Proc. 17th Ann. European Symp. on Algorithms (ESA'09)}, pp.
  444--455. Springer, 2009.
\newblock [\epfmtdoi{10.1007/978-3-642-04128-0\_40}]

\bibitem{ChekuriM09}\bibhead{ChekuriM09}
{\sc Chandra Chekuri and Benjamin Moseley}: Online scheduling to minimize the
  maximum delay factor.
\newblock In {\em Proc. 20th Ann. ACM-SIAM Symp. on Discrete Algorithms
  (SODA'09)}. SIAM, 2009.
\newblock [\epfmt{acm}{1496891}]

\bibitem{ChrobakDJKK06}\bibhead{ChrobakDJKK06}
{\sc Marek Chrobak, Christoph D\"{u}rr, Wojciech Jawor, {\L}ukasz Kowalik, and
  Maciej Kurowski}: A note on scheduling equal-length jobs to maximize
  throughput.
\newblock {\em J. of Scheduling}, 9(1):71--73, 2006.
\newblock [\epfmtdoi{10.1007/s10951-006-5595-4}]

\bibitem{DanSS96}\bibhead{DanSS96}
{\sc Asit Dan, Dinkar Sitaram, and Perwez Shahabuddin}: Dynamic batching
  policies for an on-demand video server.
\newblock {\em Multimedia Syst.}, 4(3):112--121, 1996.
\newblock [\epfmtdoi{10.1007/s005300050016}]

\bibitem{Deb84}\bibhead{Deb84}
{\sc Rajat~K. Deb}: Optimal control of bulk queues with compound {P}oisson
  arrivals and batch service.
\newblock {\em Opsearch.}, 21:227--245, 1984.

\bibitem{DebS73}\bibhead{DebS73}
{\sc Rajat~K. Deb and Richard~F. Serfozo}: Optimal control of batch service
  queues.
\newblock {\em Adv. Appl. Prob.}, 5:340--361, 1973.
\newblock [\epfmtdoi{10.2307/1426040}]

\bibitem{EdmondsP03}\bibhead{EdmondsP03}
{\sc Jeff Edmonds and Kirk Pruhs}: Multicast pull scheduling: When fairness is
  fine.
\newblock {\em Algorithmica}, 36(3):315--330, 2003.
\newblock [\epfmtdoi{10.1007/s00453-003-1018-5}]

\bibitem{EdmondsP04}\bibhead{EdmondsP04}
{\sc Jeff Edmonds and Kirk Pruhs}: A maiden analysis of longest wait first.
\newblock {\em ACM Trans. Algorithms}, 1(1):14--32, 2005.
\newblock [\epfmtdoi{10.1145/1077464.1077467}]

\bibitem{EdmondsP09}\bibhead{EdmondsP09}
{\sc Jeff Edmonds and Kirk Pruhs}: Scalably scheduling processes with arbitrary
  speedup curves.
\newblock In {\em Proc. 20th Ann. ACM-SIAM Symp. on Discrete Algorithms
  (SODA'09)}, pp. 685--692. SIAM, 2009.
\newblock [\epfmtdoi{10.1145/1496770.1496845}]

\bibitem{ErlebachH02}\bibhead{ErlebachH02}
{\sc Thomas Erlebach and Alexander Hall}: {NP}-hardness of broadcast scheduling
  and inapproximability of single-source unsplittable min-cost flow.
\newblock {\em Journal of Scheduling}, 7:223--241, 2004.
\newblock Preliminary version in
  \href{https://dl.acm.org/citation.cfm?id=545381.545405}{SODA'02}.
\newblock [\epfmtdoi{10.1023/B:JOSH.0000019682.75022.96}]

\bibitem{GandhiKKW04}\bibhead{GandhiKKW04}
{\sc Rajiv Gandhi, Samir Khuller, Yoo-Ah Kim, and Yung-Chun~(Justin) Wan}:
  Algorithms for minimizing response time in broadcast scheduling.
\newblock {\em Algorithmica}, 38(4):597--608, 2004.
\newblock Preliminary version in
  \href{https://dl.acm.org/citation.cfm?id=660079}{IPCO'02}.
\newblock [\epfmtdoi{10.1007/s00453-003-1058-x}]

\bibitem{GandhiKPS06}\bibhead{GandhiKPS06}
{\sc Rajiv Gandhi, Samir Khuller, Srinivasan Parthasarathy, and Aravind
  Srinivasan}: Dependent rounding and its applications to approximation
  algorithms.
\newblock {\em J. ACM}, 53(3):324--360, 2006.
\newblock [\epfmtdoi{10.1145/1147954.1147956}]

\bibitem{ImM10}\bibhead{ImM10}
{\sc Sungjin Im and Benjamin Moseley}: An online scalable algorithm for average
  flow time in broadcast scheduling.
\newblock In {\em Proc. 21st Ann. ACM-SIAM Symp. on Discrete Algorithms
  (SODA'10)}, pp. 1322--1333. SIAM, 2010.
\newblock [\epfmt{acm}{1873601.1873708}]

\bibitem{KalyanasundaramP95}\bibhead{KalyanasundaramP95}
{\sc Bala Kalyanasundaram and Kirk Pruhs}: Speed is as powerful as
  clairvoyance.
\newblock {\em J. ACM}, 47(4):617--643, 2000.
\newblock Preliminary version in
  \href{http://dx.doi.org/10.1109/SFCS.1995.492478}{FOCS'95}.
\newblock [\epfmtdoi{10.1145/347476.347479}]

\bibitem{KalyanasundaramPV00}\bibhead{KalyanasundaramPV00}
{\sc Bala Kalyanasundaram, Kirk~R. Pruhs, and Mahendran Velauthapillai}:
  Scheduling broadcasts in wireless networks.
\newblock {\em J. Scheduling}, 4(6):339--354, 2001.
\newblock Preliminary version in
  \href{http://dx.doi.org/10.1007/3-540-45253-2\_27}{ESA'00}.
\newblock [\epfmtdoi{10.1002/jos.87}]

\bibitem{Karger99}\bibhead{Karger99}
{\sc David Karger, Cliff Stein, and Joel Wein}: Scheduling algorithms.
\newblock In {\sc Mikhail~J. Atallah}, editor, {\em Handbook on Algorithms and
  Theory of Computation}, chapter~35. CRC Press, 1999.
\newblock [\epfmtdoi{10.1201/9781420049503-c36}]

\bibitem{KhullerK04}\bibhead{KhullerK04}
{\sc Samir Khuller and Yoo-Ah Kim}: Equivalence of two linear programming
  relaxations for broadcast scheduling.
\newblock {\em Oper. Res. Lett.}, 32(5):473--478, 2004.
\newblock [\epfmtdoi{10.1016/j.orl.2003.11.012}]

\bibitem{Kimc04}\bibhead{Kimc04}
{\sc Jae-Hoon Kim and Kyung-Yong Chwa}: Scheduling broadcasts with deadlines.
\newblock {\em Theoret. Comput. Sci.}, 325(3):479--488, 2004.
\newblock Preliminary version in
  \href{http://dx.doi.org/10.1007/3-540-45071-8_42}{COCOON'03}.
\newblock [\epfmtdoi{10.1016/j.tcs.2004.02.047}]

\bibitem{RealtimeHandbook}\bibhead{RealtimeHandbook}
{\sc Insup Lee, Joseph Y-T. Leung, and Sang~H. Son}, editors.
\newblock {\em Handbook of Real-Time and Embedded Systems}.
\newblock CRC Press, 2007.
\newblock [\epfmtdoi{10.1201/9781420011746}]

\bibitem{Pruhs07}\bibhead{Pruhs07}
{\sc Kirk Pruhs}: Competitive online scheduling for server systems.
\newblock {\em SIGMETRICS Perform. Eval. Rev.}, 34(4):52--58, 2007.
\newblock [\epfmtdoi{10.1145/1243401.1243411}]

\bibitem{PruhsST}\bibhead{PruhsST}
{\sc Kirk Pruhs, Ji\v{r}\'{i} Sgall, and Eric Torng}: Online scheduling.
\newblock In {\sc Joseph Y-T. Leung}, editor, {\em Handbook of Scheduling:
  Algorithms, Models, and Performance Analysis}. CRC Press, 2004.

\bibitem{PruhsU03}\bibhead{PruhsU03}
{\sc Kirk~R. Pruhs and Patchrawat Uthaisombut}: A comparison of multicast pull
  models.
\newblock {\em Algorithmica}, 42(3-4):289--307, 2005.
\newblock Preliminary version in
  \href{http://dx.doi.org/10.1007/3-540-45749-6\_70}{ESA'02}.
\newblock [\epfmtdoi{10.1007/s00453-005-1170-1}]

\bibitem{Sgall98}\bibhead{Sgall98}
{\sc Ji\v{r}\'{i} Sgall}: On-line scheduling.
\newblock In {\em Developments from a June 1996 seminar on Online algorithms},
  pp. 196--231. Springer-Verlag, 1998.
\newblock [\epfmt{acm}{723918}]

\bibitem{Uzsoy95}\bibhead{Uzsoy95}
{\sc Reha Uzsoy}: Scheduling batch processing machines with incompatible job
  families.
\newblock {\em International Journal of Production Research},
  33(10):2685--2708, 1995.
\newblock [\epfmtdoi{10.1080/00207549508904839}]

\bibitem{ZhengFCCPW06}\bibhead{ZhengFCCPW06}
{\sc Feifeng Zheng, Stanley P.~Y. Fung, Wun-Tat Chan, Francis Y.~L. Chin,
  Chung~Keung Poon, and Prudence W.~H. Wong}: Improved on-line broadcast
  scheduling with deadlines.
\newblock In {\em Proc. 12th Ann. Internat. Computing and Combinatorics Conf.
  (COCOON'06)}, pp. 320--329. Springer, 2006.
\newblock [\epfmtdoi{10.1007/11809678\_34}]

\end{thebibliography}
\endgroup
\end{document}
